\section{Introduction}\label{sec:intro}

Generalized Reactivity of Rank~1 (\GRone) is a fragment of linear
temporal logic (LTL) that naturally captures a broad class of
reactive systems in which a system must maintain ongoing obligations
provided the environment satisfies its own fairness
conditions~\cite{PitermanPS06}.
Several real-world systems, from hardware bus
protocols~\cite{BloemGJPPW07} to robotic motion
planning~\cite{KressGazitFP09} and autonomous
controllers~\cite{MaozR21}, are expressed as \GRone specifications.
\GRone is particularly attractive because of its tractability for
reactive synthesis: while synthesis from general LTL specifications
is doubly exponential in the length of the
formula~\cite{PnueliR89}, \GRone synthesis is polynomial in the
size of the state space~\cite{PitermanPS06}.
Tools such as Slugs~\cite{EhlersR16} and Spectra~\cite{MaozR21}
efficiently synthesize correct-by-construction controllers from
\GRone specifications.
The AMBA AHB bus arbiter, a widely cited example of synthesis
applied to an industrial hardware protocol, was specified and
synthesized as a \GRone problem~\cite{BloemGJPPW07}.

However, writing correct \GRone specifications remains a manual and
error-prone process.
Moreover, in many practical settings, a formal model may not be
available.
Consider a blackbox hardware controller or a legacy robotic system
for which only execution traces are available.
If we can obtain positive traces, representing correct system
behaviour, and negative traces, representing undesirable behaviour
(for instance, from test logs with known-good and known-bad
executions, or from traces flagged by runtime monitors), we can
mine a \GRone specification consistent with the observed behaviour.
Such a mined specification not only reveals the underlying safety
and liveness structure of the system, but can be fed directly into
existing \GRone synthesis tools to produce a correct-by-construction
controller.
This enables a pipeline from observed behaviour to formally verified
implementation, with an explicit model of the system.

Existing specification mining
tools~\cite{NeiderG18,LemieuxDB15,RahaRFV22} learn general LTL
formulas from positive and negative example traces.
However, these tools mine monolithic LTL formulas without any
distinction between environment assumptions and system guarantees.
Without this decomposition, any subsequent synthesis step must treat
all inputs as adversarial, frequently rendering the mined
specification unrealisable even when a realisable \GRone formulation
of the same behaviour exists.
Recently, ATLAS~\cite{ZhangCGD25} generalises LTL learning by
allowing users to specify structural constraints on the formula to
be learned, reducing the problem to MaxSAT via an encoding in
first-order relational logic.
One can encode the syntactic shape of a \GRone specification as an
ATLAS constraint; however, this approach treats \GRone structure as a
filter on a general LTL search space rather than exploiting it
algorithmically.
The underlying MaxSAT encoding retains the full generality of LTL
semantics over lasso traces, and scalability prohibitively degrades
as formula size and the number of traces grow.

In this paper, we present \toolname, a SAT-based algorithm that
mines \GRone specifications directly from lasso-shaped traces.
\toolname fixes the temporal skeleton of the \GRone template---the
$\always$, $\always\eventually$, and initial-condition wrappers are
given---and encodes only the unknown propositional content of each
component as an independent syntax DAG in a single SAT instance,
with hardcoded temporal semantics linking DAG evaluations to lasso
positions and consistency constraints enforcing the
assume--guarantee classification of traces.
An environment--system variable partition restricts which
propositions may appear in each component, ensuring that every
mined formula respects the assume--guarantee boundary and is
directly suitable for synthesis.
\toolname further exploits the structure of the search space via
incremental SAT solving: component DAG constraints are encoded
once and shared across template configurations using push/pop,
so that only lightweight consistency constraints are added and
retracted as the algorithm enumerates over the number of justice
and guarantee conditions.

We evaluate \toolname on 60 boolean \GRone specifications from the
Spectra benchmark suite~\cite{MaozR21}, covering AMBA bus arbiters,
generalized buffers, and parameterised arbiters with 16--218
variables.
Compared to ATLAS in both unconstrained (ATLAS[LTL]) and
\GRone-constrained (ATLAS[GR1]) configurations, \toolname solves
all benchmarks up to 34~variables within a 10-minute timeout, while
ATLAS[GR1] solves only 8/22 and ATLAS[LTL] 13/22 on the
$\leq$28-variable subset.
The formulas mined by ATLAS[LTL] are predominantly trivial
($\eventually x_i$, tautologies), yielding 1-state controllers;
\toolname produces structured \GRone specifications whose
synthesised controllers require 6--73~BDD nodes.

Section~\ref{sec:prelims} reviews \GRone and lasso traces.
Section~\ref{sec:algo} presents the SAT encoding, template
enumeration, and incremental solving.
Section~\ref{sec:eval} evaluates \toolname against ATLAS on the
Spectra benchmarks.
Section~\ref{sec:related} discusses related work, and
Section~\ref{sec:conclusion} concludes.
